\centerline{\Huge \bf  Экзаменационная работа по алгебре 8 класс}
\centerline{\Huge \bf курс углубленного изучения математики}
\centerline{\Large \bf вариант Q}

\begin{flushright}
    \scriptsize{Логика может привести Вас от пункта А к пункту Б,\\ а воображение — куда угодно.\\ Альберт Эйнштейн.}
\end{flushright}

\problem{1}
Упростите 
\[f(x) = \dfrac{1}{2} + \dfrac{1 - 4x^2}{x^3 - x^2 - x + 1} : \left( \dfrac{x}{x^2 - 2x + 1} - \dfrac{1}{1-x} \cdot \dfrac{1}{x + 1} - \dfrac{2}{x+1} \right) + x + \dfrac{1}{x^{- 1}}\]

Построить график $y = f(x)$ с учетом проведенных упрощений.

\problem{2}
Сравните число $\frac{(\sqrt{3} - 1)^3}{3\sqrt{5} - 5} - \left( \sqrt{83 - 18\sqrt{2}} + \sqrt{2} \right)$ с меньшим корнем уравнения \\ $x^2 - \left( 6 - \sqrt{2} \right)x + 8 - 2\sqrt{2} = 0$.

\problem{3}
Найдите все значения $a$ и $b$, при которых многочлен $P(x) = 2x^3 - x^2 + ax + b$ делится на $Q(x) = x^2 - 1$.

\problem{4}
Решите уравнения:

а) $82x^2(x + 5)^2 = 81 + x^4(x + 5)^4$

б) $-0.7x = \dfrac{6x}{x^3 - 4} + \dfrac{x}{x^3 + 3}$

\problem{5}
При каких значениях параметра $a$ уравнение $\frac{ax^2 - (3a + 2)x + 8}{x^2 - x - 2} = 0$ имеет единственный корень?

\problem{6}
а) Построить график функции $y = 
\begin{cases}
    \sqrt{(12 - 6 \sqrt{3})x^2} + |2x \cos 30^\circ| + x, & \text{при } x \leq 3 \\
    \frac{x - 3}{9 - x^2}, & \text{при } x > 3
\end{cases}$

б) При каких значениях $m$ прямая $y=-m$ пересечет график функции в двух точках?


\problem{7}
Алина забыла продержурить после урока и на доске остался график функции  $y = \dfrac{k}{x}$  и пять прямых, параллельных прямой  $y = kx \ (k \neq 0)$. Найдите произведение абсцисс всех точек пересечения прямых с гиперболой.

\begin{flushright}
    \textit{\small{От ЭлМШ}}
\end{flushright}

\problem{8}
Химическая лаборатория производит эксперимент со слитками. В первом слитке содержится $12$ кг золота, а во втором $6$ кг. Процент содержания золота во втором слитке на $40$ больше, чем процент содержания золота в первом. Оба слитка сплавили и получили слиток, содержащий $36$\% золота. Найти процентное содержание золота в первом слитке.

\newpage

\hspace{3mm}

\problem{9}
У Алтаны, Баира и Саши было по одинаковому тортику. Алтана разделила свой торт на $a$ равных частей и съела один кусок, Баир поделил на $b$ равных частей, съел один кусок и наконец Саша поделила свой торт на $c$ равных частей и тоже съела кусок. Так получилось, что они съели половину от одного торта. Найдите всевозможные варианты чисел $a, b$ и $c$.

\begin{flushright}
    \textit{\small{От ЭлМШ}}
\end{flushright}
